\documentclass[11pt]{article}

\usepackage[margin=1in]{geometry}
\usepackage{amsmath,amssymb,amsfonts,amsthm}
\usepackage{mathtools}
\usepackage{mathrsfs}
\usepackage{booktabs}
\usepackage{enumitem}
\usepackage[T1]{fontenc}
\usepackage[utf8]{inputenc}
\usepackage{lmodern}
\usepackage{xcolor}
\usepackage{hyperref}

\hypersetup{
  colorlinks=true,
  linkcolor=blue!50!black,
  citecolor=blue!50!black,
  urlcolor=blue!50!black
}

\newtheorem{theorem}{Theorem}[section]
\newtheorem{lemma}[theorem]{Lemma}
\newtheorem{proposition}[theorem]{Proposition}
\newtheorem{corollary}[theorem]{Corollary}
\newtheorem{example}[theorem]{Example}
\theoremstyle{definition}
\newtheorem{definition}[theorem]{Definition}
\theoremstyle{remark}
\newtheorem{remark}[theorem]{Remark}

\newcommand{\Rop}{\mathscr{R}}
\newcommand{\Phiop}{\Phi}
\newcommand{\Qop}{Q}
\newcommand{\sgn}{\operatorname{sgn}}

\title{A Local Cubic Refinement Law for Proportional-Subtended Angle Division}
\author{Wayne Baker\\Independent Researcher\\Goulds, Newfoundland and Labrador, Canada}
\date{July 2, 2026}

\begin{document}
\maketitle

\begin{abstract}
We prove a local cubic refinement law for approximate angle division generated
by proportional-subtended chord transfer.  For positive integers $a$ and $b$,
consider the seed operator
\[
D_{a,b}(\theta)
=
2\arcsin\!\left(
\frac{a}{b}\sin\frac{\theta}{2a}
\right),
\]
whose intended target is $\theta/b$.  On the local principal branch,
\[
D_{a,b}(\theta)-\frac{\theta}{b}
=
\frac{a^2-b^2}{24a^2b^3}\theta^3+O(\theta^5).
\]
For a signed error $e$, define
\[
\Phiop_{a,b}(x)
=
2\arcsin\!\left(
\frac{a}{b}\sin\frac{x}{2}
\right),
\qquad
\Rop_{a,b}(e)
=
e-
\Phiop_{a,b}\!\left(\frac{b}{a}e\right).
\]
Then, for every nonidentity branch $a\ne b$,
\[
\Rop_{a,b}(e)
=
-
\frac{a^2-b^2}{24a^2}e^3+O(e^5).
\]
The defect scaling $b/a$ is the unique constant scaling that cancels the
linear term in the one-transfer correction class, so the cubic term is sharp
within that class.  Repeated residual correction triples the local error
order, producing the hierarchy
\[
3,\quad 9,\quad 27,\quad 81,\ldots .
\]
The theorem is local and asymptotic; it is not a global angle-division theorem.
A direct finite compass-and-straightedge realization is proved for the
repeated-bisection subclass $a=2^m$.  A final two-sample chord calculation
shows how fifth-order residuals can arise once a second chord sample is
admitted, indicating a natural higher-order extension beyond the one-transfer
operator studied here.
\end{abstract}

\section*{Overview}
\label{sec:overview}

This paper isolates a local mechanism behind a family of approximate angle
division operators.  The mechanism is simple but rigid: a proportional chord
transfer has a linear part and a cubic defect; when the signed error is scaled
by $b/a$, the linear part cancels exactly.  The first surviving term is cubic.

The claims are deliberately separated.

\begin{enumerate}[label=\arabic*.]
\item The seed operator $D_{a,b}$ has an explicit cubic local error.

\item The residual operator $\Rop_{a,b}$ has an explicit cubic local error law.

\item The scaling $b/a$ is the unique constant scaling that cancels the linear
term in the one-transfer residual class.

\item Therefore, repeated residual correction triples the local order of the
remaining error.

\item A direct Euclidean realization is proved for the repeated-bisection
subclass $a=2^m$.

\item A two-sample chord calculation demonstrates that fifth-order residuals
are possible once the construction is allowed an additional chord sample.
\end{enumerate}

The analytic theorem and the finite construction theorem are related but not
identical.  The local error law holds for every positive integer pair $a\ne b$
on the local principal branch.  The direct construction statement is narrower:
it gives a sufficient compass-and-straightedge realization when the required
defect division is obtained by repeated bisection.

\section{Introduction}
\label{sec:introduction}

Classical exact angle division is constrained by Euclidean constructibility.
The best-known example is the impossibility of trisecting an arbitrary angle
by straightedge and compass alone.  The present paper addresses a different
question: how certain finite geometric approximation operators behave locally
under residual correction.

The construction begins with a proportional-subtended seed.  A chord
subtending an auxiliary angle at one radius is transferred proportionally to a
second radius, and the new subtended angle is used as an approximation to a
desired divided angle.  The same transfer operation is then applied to a
signed defect in order to correct the approximation.

The motivating $4:3$ instance is Baker's approximate-trisection
construction. Rostamian analyzed this special case, deriving its local cubic
seed error and the ninth-order error produced by one residual improvement
\cite{RostamianBaker}. The present paper begins from that branch but addresses
a more general problem: it derives the local law for arbitrary
proportional-subtended parameters $(a,b)$, proves that $b/a$ is the unique
constant defect scaling that cancels the linear term in the one-transfer
class, and establishes the resulting iterated cubic hierarchy.

The central point is not merely that the resulting error is small.  The local
error map can be computed explicitly.  Its linear term cancels identically
because the defect is scaled by $b/a$, and the first nonzero term is cubic:
\[
\Rop_{a,b}(e)
=
-
\frac{a^2-b^2}{24a^2}e^3+O(e^5).
\]
This makes the cubic order a structural consequence of the chord-transfer
operator rather than a numerical coincidence.

The result should be compared with, but not confused with, classical cubic
iteration schemes such as Halley--Householder methods.  This paper does not
claim a new general theory of third-order iteration.  Its contribution is
more specific: it identifies a geometrically interpretable residual map whose
third-order behavior follows from the proportional-subtended transfer law.
The nonlinearity enters through the chord-transfer operator
\[
\Phiop_{a,b}(x)
=
2\arcsin\!\left(
\frac{a}{b}\sin\frac{x}{2}
\right),
\]
not through an arbitrary numerical correction applied to a linear equation.

All asymptotic expansions in the paper are local expansions in radians.

\section{Branch Conventions and Chord-Transfer Interpretation}
\label{sec:branch-conventions}

All inverse trigonometric functions are interpreted on their principal real
branches.  The local theory is carried out near the origin, where all arcsine
arguments remain in $[-1,1]$ and the selected branches define single-valued
analytic representatives.

The basic transfer map is
\[
\Phiop_{a,b}(x)
=
2\arcsin\!\left(
\frac{a}{b}\sin\frac{x}{2}
\right).
\]
Its real-valuedness condition is
\[
\left|
\frac{a}{b}\sin\frac{x}{2}
\right|
\le 1.
\]
If $a\le b$, this inequality is automatically satisfied for every real $x$,
although the local principal-branch convention still matters.  If $a>b$, a
sufficient symmetric local condition is
\[
|x|
\le
2\arcsin\frac{b}{a}.
\]
For residual refinement the transfer argument is
\[
x=\frac{b}{a}e.
\]
Thus, when $a>b$, a convenient sufficient local condition is
\[
|e|
\le
\frac{2a}{b}\arcsin\frac{b}{a}.
\]
When $a\le b$, the arcsine input is bounded in magnitude by $a/b$, so the
real-valuedness condition is automatic; nevertheless, the theorem remains a
local statement near $e=0$.

The trigonometric notation is an analytic model of a Euclidean chord
operation.  A term $\sin(x/2)$ represents the half-chord ratio of a central
angle $x$.  A term $2\arcsin(s)$ represents the central angle whose half-chord
ratio is $s$ on the chosen branch.  In the finite construction statements,
these operations are realized by constructing chords, transferring lengths,
and recovering central angles, not by numerical evaluation of transcendental
functions.

\section{The Proportional-Subtended Seed}
\label{sec:seed}

\begin{definition}[Proportional-subtended seed]
Let $a$ and $b$ be positive integers.  Define
\[
D_{a,b}(\theta)
=
2\arcsin\!\left(
\frac{a}{b}\sin\frac{\theta}{2a}
\right).
\]
The intended target is $\theta/b$, and the signed seed error is
\[
e^{(a,b)}_0(\theta)
=
D_{a,b}(\theta)-\frac{\theta}{b}.
\]
\end{definition}

\begin{lemma}[Seed expansion]
\label{lem:seed-expansion}
For positive integers $a$ and $b$,
\[
D_{a,b}(\theta)
=
\frac{\theta}{b}
+
\frac{a^2-b^2}{24a^2b^3}\theta^3
+
O(\theta^5).
\]
Consequently,
\[
e^{(a,b)}_0(\theta)
=
\frac{a^2-b^2}{24a^2b^3}\theta^3
+
O(\theta^5).
\]
\end{lemma}

\begin{proof}
Use the local expansions
\[
\sin x=x-\frac{x^3}{6}+O(x^5),
\qquad
\arcsin y=y+\frac{y^3}{6}+O(y^5).
\]
First,
\[
\sin\frac{\theta}{2a}
=
\frac{\theta}{2a}
-
\frac{\theta^3}{48a^3}
+
O(\theta^5).
\]
Multiplying by $a/b$ gives
\[
\frac{a}{b}\sin\frac{\theta}{2a}
=
\frac{\theta}{2b}
-
\frac{\theta^3}{48a^2b}
+
O(\theta^5).
\]
Applying the arcsine expansion,
\[
\arcsin\!\left(
\frac{a}{b}\sin\frac{\theta}{2a}
\right)
=
\frac{\theta}{2b}
-
\frac{\theta^3}{48a^2b}
+
\frac{\theta^3}{48b^3}
+
O(\theta^5).
\]
After multiplication by $2$,
\[
D_{a,b}(\theta)
=
\frac{\theta}{b}
+
\left(
-
\frac{1}{24a^2b}
+
\frac{1}{24b^3}
\right)\theta^3
+
O(\theta^5).
\]
Since
\[
-
\frac{1}{24a^2b}
+
\frac{1}{24b^3}
=
\frac{a^2-b^2}{24a^2b^3},
\]
the result follows.
\end{proof}

\begin{remark}[Overshoot, undershoot, and the identity branch]
The sign of the leading seed error is determined by $a^2-b^2$.  If $a>b$,
the seed locally overshoots $\theta/b$.  If $a<b$, the seed locally
undershoots $\theta/b$.  The identity branch $a=b$ is exceptional: on the
local principal branch,
\[
D_{a,a}(\theta)=\frac{\theta}{a},
\]
so the seed error vanishes identically.
\end{remark}

\section{The Transfer Map and the Residual Law}
\label{sec:residual-law}

\begin{definition}[Transfer map]
For positive integers $a$ and $b$, define
\[
\Phiop_{a,b}(x)
=
2\arcsin\!\left(
\frac{a}{b}\sin\frac{x}{2}
\right).
\]
\end{definition}

\begin{lemma}[Transfer-map expansion]
\label{lem:transfer-expansion}
The transfer map satisfies
\[
\Phiop_{a,b}(x)
=
\frac{a}{b}x
+
\frac{a(a^2-b^2)}{24b^3}x^3
+
O(x^5).
\]
It is odd wherever it is real-valued:
\[
\Phiop_{a,b}(-x)=-\Phiop_{a,b}(x).
\]
\end{lemma}

\begin{proof}
Expanding directly,
\[
\frac{a}{b}\sin\frac{x}{2}
=
\frac{a}{2b}x
-
\frac{a}{48b}x^3
+
O(x^5),
\]
and hence
\[
\arcsin\!\left(
\frac{a}{b}\sin\frac{x}{2}
\right)
=
\frac{a}{2b}x
-
\frac{a}{48b}x^3
+
\frac{a^3}{48b^3}x^3
+
O(x^5).
\]
Multiplication by $2$ gives
\[
\Phiop_{a,b}(x)
=
\frac{a}{b}x
+
\frac{a(a^2-b^2)}{24b^3}x^3
+
O(x^5).
\]
Oddness follows because $\sin(x/2)$ and $\arcsin x$ are odd on their local
principal real branches.
\end{proof}

\begin{definition}[Signed-defect residual operator]
For a signed error $e$, define
\[
\Rop_{a,b}(e)
=
e-
\Phiop_{a,b}\!\left(\frac{b}{a}e\right).
\]
\end{definition}

\begin{theorem}[Universal proportional-subtended cubic error law]
\label{thm:universal-law}
Let $a$ and $b$ be positive integers with $a\ne b$.  Then
\[
\Rop_{a,b}(e)
=
-
\frac{a^2-b^2}{24a^2}e^3
+
O(e^5).
\]
Thus the linear term cancels identically, and the first nonzero term is cubic.
\end{theorem}

\begin{proof}
Insert $x=(b/a)e$ into Lemma~\ref{lem:transfer-expansion}:
\[
\Phiop_{a,b}\!\left(\frac{b}{a}e\right)
=
\frac{a}{b}\left(\frac{b}{a}e\right)
+
\frac{a(a^2-b^2)}{24b^3}
\left(\frac{b}{a}e\right)^3
+
O(e^5).
\]
Therefore
\[
\Phiop_{a,b}\!\left(\frac{b}{a}e\right)
=
e+
\frac{a^2-b^2}{24a^2}e^3
+
O(e^5).
\]
Subtracting this expression from $e$ gives
\[
\Rop_{a,b}(e)
=
-
\frac{a^2-b^2}{24a^2}e^3
+
O(e^5).
\]
Because $a\ne b$, the cubic coefficient is nonzero.
\end{proof}

\section{Uniqueness of the Scaling and One-Transfer Sharpness}
\label{sec:sharpness}

The defect scaling $b/a$ is not a heuristic tuning parameter.  It is the
unique constant scaling that cancels the linear term of the proportional-
subtended transfer map.  This makes the cubic law sharp within the correction
class that uses one transfer map applied to one constant multiple of the signed
defect.

Let
\[
r=\frac{a}{b}.
\]
Then
\[
\Phiop_{a,b}(x)
=
2\arcsin\!\left(r\sin\frac{x}{2}\right)
=
rx+
\frac{r(r^2-1)}{24}x^3
+
O(x^5).
\]
For a constant $\lambda$, consider
\[
\Rop^{(a,b)}_{\lambda}(e)
=
e-
\Phiop_{a,b}(\lambda e).
\]
Then
\[
\Rop^{(a,b)}_{\lambda}(e)
=
(1-r\lambda)e
-
\frac{r(r^2-1)\lambda^3}{24}e^3
+
O(e^5).
\]
The linear term vanishes if and only if
\[
1-r\lambda=0,
\]
or equivalently
\[
\lambda=\frac{1}{r}=\frac{b}{a}.
\]
For this unique value,
\[
\Rop^{(a,b)}_{b/a}(e)
=
-
\frac{r^2-1}{24r^2}e^3
+
O(e^5)
=
-
\frac{a^2-b^2}{24a^2}e^3
+
O(e^5).
\]
For every nonidentity branch $a\ne b$, this cubic coefficient is nonzero.
Therefore no one-transfer correction of the form
\[
e-
\Phiop_{a,b}(\lambda e)
\]
can have local order greater than three.

This sharpness is a ceiling only for the one-transfer, constant-scaled defect
class.  It is not an obstruction to higher local order once additional chord
samples or other constructible degrees of freedom are introduced.

\section{The Cubic Cascade and Leading Coefficients}
\label{sec:cubic-cascade}

Define the seed approximation by
\[
X^{(a,b)}_0(\theta)=D_{a,b}(\theta),
\]
and define recursively
\[
X^{(a,b)}_{k+1}(\theta)
=
\frac{\theta}{b}
+
\Rop_{a,b}\!\left(
X^{(a,b)}_k(\theta)-\frac{\theta}{b}
\right).
\]
Equivalently, if
\[
e^{(a,b)}_k(\theta)
=
X^{(a,b)}_k(\theta)-\frac{\theta}{b},
\]
then
\[
e^{(a,b)}_{k+1}(\theta)
=
\Rop_{a,b}\!\left(e^{(a,b)}_k(\theta)\right).
\]

\begin{theorem}[General proportional-subtended cubic cascade]
\label{thm:general-cubic-cascade}
Suppose that, for some $k\ge 0$,
\[
e^{(a,b)}_k(\theta)
=
c^{(a,b)}_k\theta^{q_k}
+
O(\theta^{q_k+2}),
\qquad
c^{(a,b)}_k\ne 0.
\]
Then
\[
e^{(a,b)}_{k+1}(\theta)
=
-
\frac{a^2-b^2}{24a^2}
\left(c^{(a,b)}_k\right)^3
\theta^{3q_k}
+
O(\theta^{3q_k+2}).
\]
Consequently,
\[
q_{k+1}=3q_k.
\]
Since the seed has $q_0=3$, one obtains
\[
q_k=3^{k+1}.
\]
Thus the local error orders are
\[
3,\quad 9,\quad 27,\quad 81,\ldots .
\]
\end{theorem}

\begin{proof}
By Theorem~\ref{thm:universal-law},
\[
\Rop_{a,b}(e)
=
-
\frac{a^2-b^2}{24a^2}e^3
+
O(e^5).
\]
Thus
\[
e^{(a,b)}_{k+1}(\theta)
=
-
\frac{a^2-b^2}{24a^2}
\left(e^{(a,b)}_k(\theta)\right)^3
+
O\!\left(\left(e^{(a,b)}_k(\theta)\right)^5\right).
\]
Now assume
\[
e^{(a,b)}_k(\theta)
=
c^{(a,b)}_k\theta^{q_k}
+
O(\theta^{q_k+2}).
\]
Cubing gives
\[
\left(e^{(a,b)}_k(\theta)\right)^3
=
\left(c^{(a,b)}_k\right)^3\theta^{3q_k}
+
O(\theta^{3q_k+2}),
\]
because the first cross-term has order $2q_k+(q_k+2)=3q_k+2$.  Also,
\[
\left(e^{(a,b)}_k(\theta)\right)^5
=
O(\theta^{5q_k}).
\]
For the cascade generated from the seed, $q_k\ge 3$, so $5q_k\ge 3q_k+2$.
The $O(e_k^5)$ contribution is therefore absorbed into
$O(\theta^{3q_k+2})$.  This proves the displayed recurrence and hence
$q_{k+1}=3q_k$.  Since Lemma~\ref{lem:seed-expansion} gives $q_0=3$, induction
gives $q_k=3^{k+1}$.
\end{proof}

\begin{corollary}[Leading-coefficient recurrence]
\label{cor:leading-coefficients}
The leading coefficients satisfy
\[
c^{(a,b)}_{k+1}
=
-
\frac{a^2-b^2}{24a^2}
\left(c^{(a,b)}_k\right)^3,
\]
with
\[
c^{(a,b)}_0
=
\frac{a^2-b^2}{24a^2b^3}.
\]
Equivalently,
\[
c^{(a,b)}_k
=
(-1)^k
\left(
\frac{a^2-b^2}{24a^2}
\right)^{(3^k-1)/2}
\left(
\frac{a^2-b^2}{24a^2b^3}
\right)^{3^k}.
\]
\end{corollary}

\begin{proof}
The recurrence follows directly from
Theorem~\ref{thm:general-cubic-cascade}.  Iterating it gives the closed form:
the exponent of the residual coefficient is
\[
0,\quad 1,\quad 4,\quad 13,\ldots,
\]
namely $(3^k-1)/2$, while the exponent of the seed coefficient is multiplied
by $3$ at every stage, giving $3^k$.  The sign alternates once per refinement,
giving $(-1)^k$.
\end{proof}

\begin{table}[h]
\centering
\begin{tabular}{cccc}
\toprule
Branch & Target & Seed sign & Local order cascade \\
\midrule
$1:2$ & $\theta/2$ & negative & $3,9,27,81,\ldots$ \\
$2:3$ & $\theta/3$ & negative & $3,9,27,81,\ldots$ \\
$3:2$ & $\theta/2$ & positive & $3,9,27,81,\ldots$ \\
$4:3$ & $\theta/3$ & positive & $3,9,27,81,\ldots$ \\
$3:4$ & $\theta/4$ & negative & $3,9,27,81,\ldots$ \\
$5:4$ & $\theta/4$ & positive & $3,9,27,81,\ldots$ \\
\bottomrule
\end{tabular}
\caption{Representative proportional-subtended branches.  The sign records the
leading sign of the seed error.}
\label{tab:representative-branches}
\end{table}

\section{Euclidean Realization}
\label{sec:geometry}

The analytic refinement argument is
\[
\frac{b}{a}e
=
\frac{bX-\theta}{a},
\]
where $X$ is a current approximation to $\theta/b$.  The angle $bX-\theta$ can
be formed by repeated angle copying and subtraction.  The remaining operation
is division of that constructed defect angle by $a$, followed by the
proportional-subtended transfer map.

\begin{lemma}[Euclidean realization of the transfer map]
\label{lem:transfer-realization}
Let $a$ and $b$ be positive constructible radii, and let $x$ be a constructed
signed angle.  Suppose
\[
\left|
\frac{a}{b}\sin\frac{x}{2}
\right|
\le 1.
\]
Then the angle
\[
\Phiop_{a,b}(x)
=
2\arcsin\!\left(
\frac{a}{b}\sin\frac{x}{2}
\right)
\]
is constructible from $x$, $a$, and $b$.
\end{lemma}

\begin{proof}
Construct a circle of radius $a$.  On this circle, construct the chord
subtended by the signed central angle $x$.  Its directed chord length is
\[
2a\sin\frac{x}{2}.
\]
Transfer this chord length to a circle of radius $b$.  Let $\Phi$ be the signed
central angle subtended by the transferred chord on the radius-$b$ circle.
Then
\[
2b\sin\frac{\Phi}{2}
=
2a\sin\frac{x}{2}.
\]
Therefore
\[
\sin\frac{\Phi}{2}
=
\frac{a}{b}\sin\frac{x}{2},
\]
and, on the chosen local principal branch,
\[
\Phi
=
2\arcsin\!\left(
\frac{a}{b}\sin\frac{x}{2}
\right)
=
\Phiop_{a,b}(x).
\]
Thus the transfer map is realized by chord construction and chord transfer.
\end{proof}

\begin{proposition}[Sufficient repeated-bisection realization criterion]
\label{prop:bisection-realization}
Suppose $a=2^m$ for some integer $m\ge 0$.  Then the defect angle
\[
\frac{bX-\theta}{a}
\]
can be obtained from $bX-\theta$ by $m$ successive bisections.  Hence, whenever
the required copied and subtracted angles are available in the chosen geometric
setup and the transfer lies in the real branch domain, the proportional-
subtended refinement step has a direct repeated-bisection realization.
\end{proposition}

\begin{proof}
Repeated angle copying constructs $bX$.  Subtracting the input angle constructs
$bX-\theta$.  If $a=2^m$, then $m$ successive bisections construct
$(bX-\theta)/a$.  Lemma~\ref{lem:transfer-realization} then realizes
$\Phiop_{a,b}$, and the sign of the defect determines the orientation of the
final correction.
\end{proof}

\begin{example}[The $4:3$ trisection branch]
For $a=4$ and $b=3$, the target is $\theta/3$.  Let $X$ be a current
approximation to $\theta/3$, and write
\[
X=\frac{\theta}{3}+e.
\]
The constructed defect angle is
\[
3X-\theta=3e.
\]
Since $a=4$, two bisections construct
\[
\frac{3X-\theta}{4}=\frac{3}{4}e.
\]
Applying the transfer map gives
\[
\Phiop_{4,3}\!\left(\frac{3X-\theta}{4}\right)
=
2\arcsin\!\left(
\frac{4}{3}\sin\frac{3X-\theta}{8}
\right).
\]
Thus the refined approximation is
\[
X_{\mathrm{new}}
=
X-
2\arcsin\!\left(
\frac{4}{3}\sin\frac{3X-\theta}{8}
\right).
\]
This is the $4:3$ residual correction associated with Baker's trisection
construction and Rostamian's special-case error analysis
\cite{RostamianBaker}.
\end{example}

\begin{remark}[Sufficient, not complete]
The condition $a=2^m$ is a sufficient criterion for the most direct
construction based on angle copying, subtraction, and repeated bisection.  It
is not claimed to classify all possible compass-and-straightedge realizations.
\end{remark}

\begin{remark}[Analytic law versus finite construction]
The local cubic law
\[
\Rop_{a,b}(e)
=
-
\frac{a^2-b^2}{24a^2}e^3+O(e^5)
\]
holds analytically for every nonidentity branch $a\ne b$.  The direct finite
construction above is proved only for the repeated-bisection subclass
$a=2^m$.  These are distinct claims.
\end{remark}

\section{A Two-Sample Chord Correction}
\label{sec:two-sample}

The one-transfer sharpness result does not preclude higher local order.  It
only says that a correction of the form
\[
e-\Phiop_{a,b}(\lambda e)
\]
cannot exceed cubic order for $a\ne b$.  Higher order requires another degree
of freedom.  A natural geometric source is a second chord sample obtained by
bisecting the defect angle.

Let
\[
r=\frac{a}{b},
\qquad
 y=\frac{e}{r}=\frac{b}{a}e.
\]
The cubic correction uses the single chord sample
\[
\sin\frac{y}{2}.
\]
If the defect angle $y$ is bisected first, then the chord of the bisected
defect supplies a second sample,
\[
\sin\frac{y}{4}.
\]
Consider the two-sample expression
\[
A\sin\frac{y}{2}
+
B\sin\frac{y}{4},
\]
and define
\[
\Qop_r(y)
=
2\arcsin\!\left(
r\left(
A\sin\frac{y}{2}
+
B\sin\frac{y}{4}
\right)
\right).
\]
We choose $A$ and $B$ so that $\Qop_r(y)$ agrees with $ry$ through cubic order.
Using
\[
\sin\frac{y}{2}
=
\frac{y}{2}-\frac{y^3}{48}+O(y^5),
\]
and
\[
\sin\frac{y}{4}
=
\frac{y}{4}-\frac{y^3}{384}+O(y^5),
\]
the matching conditions reduce to
\[
A+\frac{B}{2}=1,
\qquad
A+\frac{B}{8}=r^2.
\]
Solving gives
\[
A=\frac{4r^2-1}{3},
\qquad
B=\frac{8(1-r^2)}{3}.
\]
With these choices,
\[
\Qop_r(y)
=
ry
-
\frac{r(r^2-1)(4r^2-1)}{7680}y^5
+
O(y^7).
\]
Consequently,
\[
e-
\Qop_r\!\left(\frac{e}{r}\right)
=
\frac{(r^2-1)(4r^2-1)}{7680r^4}e^5
+
O(e^7).
\]
Thus cubic order is sharp for the one-transfer correction, but it is not an
absolute barrier once additional chord samples are admitted.

For the $4:3$ branch,
\[
r=\frac{4}{3},
\]
so
\[
A=\frac{55}{27},
\qquad
B=-\frac{56}{27}.
\]
The corresponding two-sample correction is
\[
\Qop_{4/3}(y)
=
2\arcsin\!\left(
\frac{4}{3}
\left(
\frac{55}{27}\sin\frac{y}{2}
-
\frac{56}{27}\sin\frac{y}{4}
\right)
\right),
\]
and its residual begins at fifth order:
\[
e-
\Qop_{4/3}\!\left(\frac{3e}{4}\right)
=
\frac{77}{393216}e^5
+
O(e^7).
\]

This section is included to clarify the scope of the sharpness theorem.  The
systematic construction and analysis of higher-order chord-combination
corrections is left as a future problem.

\section{Numerical Consistency Checks}
\label{sec:numerical}

The computations in this section are consistency checks only.  They are not
used in the proof of the local cubic law, and they do not enlarge the finite
geometric-realization claim.  In particular, the computations do not establish
global convergence and do not classify which branches admit compass-and-
straightedge realizations.

For representative branches, we evaluated both the seed error
\[
D_{a,b}(\theta)-\frac{\theta}{b}
\]
and the refinement residual
\[
\Rop_{a,b}(e)
=
e-
\Phiop_{a,b}\!\left(\frac{b}{a}e\right).
\]
The predicted local laws are
\[
D_{a,b}(\theta)-\frac{\theta}{b}
=
\frac{a^2-b^2}{24a^2b^3}\theta^3+O(\theta^5),
\]
and
\[
\Rop_{a,b}(e)
=
-
\frac{a^2-b^2}{24a^2}e^3+O(e^5).
\]
Thus log--log fits should approach slope $3$ as the input scale tends to zero.

\begin{table}[h]
\centering
\begin{tabular}{cccc}
\toprule
Branch & Seed slope & Refinement slope & Consistency \\
\midrule
$1:2$ & $3.0046193339$ & $3.0006346375$ & Consistent \\
$2:3$ & $3.0028122658$ & $3.0008598848$ & Consistent \\
$3:2$ & $3.0080787730$ & $3.0010922823$ & Consistent \\
$4:3$ & $3.0035274023$ & $3.0010770871$ & Consistent \\
$3:4$ & $3.0016941563$ & $3.0009206657$ & Consistent \\
$5:4$ & $3.0019634516$ & $3.0010671119$ & Consistent \\
\bottomrule
\end{tabular}
\caption{Representative high-precision consistency checks for the local cubic
law.  The slopes are numerical log--log estimates and are included only to
confirm that direct evaluation reproduces the cubic behavior predicted by the
proof.}
\label{tab:numerical-slopes}
\end{table}

As a concrete scale example, take the $4:3$ branch and $\theta=60^\circ$.  The
seed error is approximately
\[
4.50\times 10^{-2}\text{ degrees}.
\]
After one residual correction, the error is approximately
\[
5.06\times 10^{-10}\text{ degrees},
\]
and after a second correction it is approximately
\[
7.20\times 10^{-34}\text{ degrees}.
\]
These numbers illustrate the order-tripling effect in the familiar trisection
branch.  The proof is the formal expansion of $\Rop_{a,b}$, not the numerical
table.

As a broader stress check for the $4:3$ branch,
Table~\ref{tab:large-angle-stress} records signed errors for several larger
input angles.  These computations are not a global convergence theorem; they
only show that, for these tested angles, the residual correction continues to
produce rapid error reduction.

\begin{table}[htbp]
\centering
\scriptsize
\setlength{\tabcolsep}{4pt}
\begin{tabular}{@{}cccc@{}}
\toprule
Input $\theta$ & Seed error $e_0$ & First residual $e_1$ & Second residual $e_2$ \\
\midrule
$45^\circ$  & $+1.887774001933852\times 10^{-2}$ & $-3.735696264404784\times 10^{-11}$ & $+2.894917723791054\times 10^{-37}$ \\
$135^\circ$ & $+5.416639994179353\times 10^{-1}$ & $-8.825017205888855\times 10^{-7}$  & $+3.816526069041368\times 10^{-24}$ \\
$225^\circ$ & $+2.883353443312359\times 10^{0}$  & $-1.331469483006689\times 10^{-4}$  & $+1.310735921489383\times 10^{-17}$ \\
$315^\circ$ & $+1.052789884545587\times 10^{1}$  & $-6.502746781381637\times 10^{-3}$  & $+1.526905220006051\times 10^{-12}$ \\
\bottomrule
\end{tabular}
\caption{Large-angle numerical stress check for the $4:3$ branch.  Errors are
signed errors in degrees, rounded to approximately double-precision reporting
level.  The table illustrates rapid residual reduction for the tested angles,
but it is not a global convergence theorem.}
\label{tab:large-angle-stress}
\end{table}

The larger examples show why the theorem is best interpreted as a local
residual law rather than merely as a small-$\theta$ seed expansion: once the
seed defect is sufficiently small, the residual map acts cubically on that
defect.

Finally, Table~\ref{tab:two-sample-fifth-order-check} gives a small numerical
check of the two-sample fifth-order residual in the $4:3$ branch.

\begin{table}[htbp]
\centering
\scriptsize
\setlength{\tabcolsep}{6pt}
\begin{tabular}{@{}ccc@{}}
\toprule
Input defect $e$ & Fifth-order residual & Observed scale \\
\midrule
$1^\circ$   & $+1.817122518225172\times 10^{-11}\!^\circ$ & $O(e^5)$ \\
$0.1^\circ$ & $+1.817059574800297\times 10^{-16}\!^\circ$ & $O(e^5)$ \\
$0.01^\circ$ & $+1.817058945385739\times 10^{-21}\!^\circ$ & $O(e^5)$ \\
\bottomrule
\end{tabular}
\caption{Numerical check of the two-sample fifth-order residual for the $4:3$
branch.  Reducing the input defect by a factor of $10$ reduces the residual by
approximately $10^5$, as predicted by the expansion.}
\label{tab:two-sample-fifth-order-check}
\end{table}

\section{Scope of the Result}
\label{sec:scope}

The result proved here is local, not global.  It is a theorem about the local
asymptotic behavior of an explicit proportional-subtended residual operator.
It does not assert that every branch gives a globally convergent angle-division
algorithm, nor does it assert that every analytic branch has a direct finite
compass-and-straightedge realization.

The scope can be summarized as follows.

\begin{enumerate}[label=\arabic*.]
\item The analytic seed law holds locally for every positive integer pair
$a,b$.

\item The analytic residual law holds locally for every nonidentity branch
$a\ne b$.

\item The direct Euclidean realization is proved for the repeated-bisection
subclass $a=2^m$.

\item The one-transfer sharpness theorem applies only to residual corrections
of the form $e-\Phiop_{a,b}(\lambda e)$.

\item The two-sample calculation shows that higher order is possible after an
additional chord sample is admitted, but it is not developed into a full
construction theory here.
\end{enumerate}

These distinctions are important.  They keep the paper's main theorem broad
where the analytic expansion is broad, and narrow where the finite construction
record is narrow.

\section{Conclusion}
\label{sec:conclusion}

The proportional-subtended family gives a local cubic refinement law for
approximate angle division.  The seed satisfies
\[
D_{a,b}(\theta)-\frac{\theta}{b}
=
\frac{a^2-b^2}{24a^2b^3}\theta^3+O(\theta^5),
\]
and the signed-defect residual operator satisfies
\[
\Rop_{a,b}(e)
=
-
\frac{a^2-b^2}{24a^2}e^3+O(e^5).
\]
The defect scaling $b/a$ is the unique constant scaling that cancels the
first-order defect in the one-transfer residual class.  Consequently, every
residual correction triples the local error order:
\[
3\longrightarrow 9\longrightarrow 27\longrightarrow 81\longrightarrow\cdots .
\]

For $a=2^m$, the residual correction has a direct repeated-bisection
straightedge-and-compass realization.  For general $a\ne b$, the local analytic
error law remains valid, while the finite construction status is a separate
question.  The two-sample chord calculation shows a natural route to
fifth-order residuals and suggests a further problem: develop a systematic
finite-construction theory for higher-order chord-combination corrections.

\appendix

\section{Symbolic Verification Recipe}
\label{app:symbolic}

The following symbolic computation verifies the displayed seed, residual, and
two-sample chord coefficients.  It is included only as a reproducibility aid;
the proofs above use the same formal series composition by hand.

\begin{verbatim}
import sympy as sp

theta, e, x, y = sp.symbols('theta e x y')
a, b, r = sp.symbols('a b r', positive=True)

D = 2*sp.asin((a/b)*sp.sin(theta/(2*a)))
print(sp.series(D - theta/b, theta, 0, 5))

Phi = 2*sp.asin((a/b)*sp.sin(x/2))
R = e - Phi.subs(x, (b/a)*e)
print(sp.series(R, e, 0, 5))

A = (4*r**2 - 1)/3
B = 8*(1 - r**2)/3

Q = 2*sp.asin(
    r*(A*sp.sin(y/2) + B*sp.sin(y/4))
)

print(sp.factor(sp.series(Q - r*y, y, 0, 7)))
print(sp.factor(sp.series(e - Q.subs(y, e/r), e, 0, 7)))
\end{verbatim}

\begin{thebibliography}{9}

\bibitem{Wantzel1837}
P.-L. Wantzel,
``Recherches sur les moyens de reconnaitre si un probleme de geometrie peut se resoudre avec la regle et le compas,''
\emph{Journal de Mathematiques Pures et Appliquees}, vol.~2, pp.~366--372, 1837.

\bibitem{Durer1525}
A. Durer,
\emph{Underweysung der Messung mit dem Zirckel und Richtscheyt},
Nuremberg, 1525.

\bibitem{Householder1970}
A. S. Householder,
\emph{The Numerical Treatment of a Single Nonlinear Equation},
McGraw--Hill, 1970.

\bibitem{Traub1964}
J. F. Traub,
\emph{Iterative Methods for the Solution of Equations},
Prentice--Hall, 1964.

\bibitem{StoerBulirsch2002}
J. Stoer and R. Bulirsch,
\emph{Introduction to Numerical Analysis}, 3rd ed.,
Texts in Applied Mathematics, vol.~12, Springer, New York, 2002.

\bibitem{RostamianBaker}
R. Rostamian,
``An angle trisection: construction by Wayne Baker,''
University of Maryland, Baltimore County, web note and error analysis,
\url{https://userpages.umbc.edu/~rostamia/Geometry/trisect-baker.html},
accessed July 13, 2026.

\end{thebibliography}

\end{document}
